\documentclass{article}
\usepackage{amsmath}
\usepackage{graphicx}
\usepackage{hyperref}
\usepackage{listings}
\usepackage{booktabs}
\usepackage{threeparttable}
\usepackage[
backend=biber,
style=authoryear-comp,
doi=false,isbn=false,url=false,eprint=false,maxcitenames=2,uniquename=false
]{biblatex}
%\setbeamertemplate{bibliography item}{\insertbiblabel}
\addbibresource{bibliography.bib}

\title{Benefits of Using \LaTeX\ for Academic Writing}
\author{E383}
\date{\today}

\begin{document}

\maketitle

\tableofcontents

\section{Introduction}

\LaTeX\ is a typesetting system commonly used for academic and scientific documents. It offers numerous benefits over traditional word processors, particularly for technical and academic writing.

\section{High-Quality Typesetting}

\LaTeX\ produces documents with professional, high-quality typesetting. It excels in handling complex formatting, including mathematical equations, tables, and bibliographies. 

\subsection{Mathematical Typesetting}

One of the standout features of \LaTeX\ is its ability to render complex mathematical equations beautifully. For example:

\begin{equation}
\label{eq:taylor}
f(x) \approx f(a) + f'(a)(x - a) + \frac{f''(a)}{2}(x - a)^2
\end{equation}

\begin{equation}
    y_{it}
\end{equation}

\subsection{Consistent Formatting}

\LaTeX\ uses a consistent approach to document formatting, ensuring that all elements such as headings, paragraphs, and references maintain a uniform style throughout the document. This is particularly useful for large documents like theses and dissertations.

\section{Automated Citation and Bibliography Management}

\LaTeX\ integrates well with BibTeX and other bibliography management tools, making it easy to handle citations and references. For example:
\textcite{wooldridgeIntroductoryEconometricsModern2016}

\section{Cross-Referencing}

Cross-referencing figures, tables, and sections is straightforward in \LaTeX. This makes it easy to create and manage complex documents. For instance, see Figure~\ref{fig:example}.

\begin{figure}[ht]
    \centering
    \includegraphics[width=0.5\textwidth]{example-image}
    \caption{An example image.}
    \label{fig:example}
\end{figure}

It is then easy to reference back to figure \ref{fig:example} or equation (\ref{eq:taylor}).

\input{sections/scalability}

\section{Tables and Figures}

Creating       complex tables and inserting figures is \\ straightforward in \LaTeX. Below is an example table created using \LaTeX.

\begin{table}[ht]
\centering
\caption{An Example Table}
\begin{threeparttable}
\begin{tabular}{lcc}
\hline
 & Coefficient & Standard error \\
\hline
           Constant & 0.171 & 0.141 \\
Initial level of intended funds rate & -0.021 & 0.012 \\
\hline
\end{tabular}
\begin{tablenotes}
\item \textit{Notes:} This is an example of a table created using \LaTeX.
\end{tablenotes}
\end{threeparttable}
\end{table}

\section{Code Listings}

\LaTeX\ supports including code listings with syntax highlighting, which is useful for computer science and engineering documents. Here is an example:

\begin{lstlisting}[language=Python, caption=Python example]
def hello_world():
    print("Hello, world!")
\end{lstlisting}

\section{Conclusion}

In summary, \LaTeX\ offers numerous benefits for academic writing, including high-quality typesetting, automated citation management, and support for complex documents. Its stability, platform independence, and customization options make it an excellent choice for researchers and academics.


\begin{enumerate}
    \item Bullet 1 \\
    Bullet 1.2
    \item Bullet 2
\end{enumerate}

\begin{equation}
    \frac{a}{b}
\end{equation}

in my text i want $\frac{a}{b}$

\printbibliography

\end{document}
